\section{Related Work}
\label{sec:related}

\subsection{Cybersickness Detection and Its Evaluation}
Questionnaires such as the Simulator Sickness Questionnaire~\cite{kennedy1993ssq} measure discomfort after exposure and cannot capture onset during immersion~\cite{caserman2021cybersickness}, which motivates sensor-driven detection. Dennison et al.~\cite{dennison2016physio} related physiological signals to cybersickness reports; Islam et al. detected and forecast severity from heart rate, breathing, heart-rate variability and skin conductance~\cite{islam2020automatic} and later fused eye- and head-tracking data~\cite{islam2021fusion}; Liu et al. used EEG with a CNN-ECA-LSTM network~\cite{liu2024cnneca}; Yalcin et al. combined bidirectional LSTMs with GAN-based augmentation of consumer-grade physiological data~\cite{yalcin2024automatic}; and Jeong and Han applied transformers to multimodal time series~\cite{jeong2024precyse}. Explainable models have been used to identify dominant features and to shrink detectors for deployment~\cite{kundu2022truvr,kundu2023litevr,kundu2023vrlens}. A systematic review notes that most detectors classify the current state rather than anticipating onset, and that evaluation protocols vary widely~\cite{yang2022mlreview}.

The gameplay data used here were collected by Porcino et al.~\cite{porcino2021methodology,porcino2019uffcsdata}, who used symbolic machine learning to identify likely causes of discomfort such as rotation and acceleration~\cite{porcino2022identifying}. The same gameplay data were used by TruVR, which reported \pct{94.10} accuracy with explainable boosting machines trained on \pct{70} of the samples and tested on the remaining \pct{30}~\cite{kundu2022truvr}, and in work on attacks against cybersickness detectors~\cite{kundu2025securing}. Splitting samples rather than participants has been used in this literature~\cite{kundu2022truvr}; with 1-Hz samples and slowly changing labels, it allows a model to recognize sessions it has already seen~\cite{kapoor2023leakage,kaufman2012leakage,roberts2017cv}. We quantify this effect on the gameplay data in Section~\ref{sec:leak}; we do not re-evaluate the cited studies, whose exact pipelines are not available to us.

\subsection{Adversarial Threats to VR Users and Detectors}
VR mediates perception and balance and therefore exposes an attack surface distinct from conventional software~\cite{odeleye2023virtually,giaretta2025security}. Tseng et al. showed that perceptual manipulations such as redirected walking can be repurposed to cause physical harm~\cite{tseng2022dark}, and Cheng et al. studied user reactions to perceptual-manipulation attacks in mixed reality~\cite{cheng2023perceptual}. Odeleye et al. described frame-rate manipulation attacks that can induce VR sickness and detected them with an isolation forest~\cite{odeleye2021framerate}. Detectors themselves can be attacked: Kundu et al. crafted adversarial input perturbations that stop deep cybersickness detectors from triggering mitigation, evaluated them on two public datasets including the gameplay data used here, and proposed an XAI-guided attack detector~\cite{kundu2025securing}; Ahmed et al. built an open-source testbed in which gradient-based attacks (MI-FGSM, PGD, C\&W) degrade deep detectors that drive a field-of-view mitigation in closed loop~\cite{ahmed2025adversarialvr}. These works perturb detector inputs by optimization; the stress tests in this paper instead apply bounded, physically motivated manipulations to the logged telemetry channels (Section~\ref{sec:threat}), and neither approach measures what a manipulation would do to a participant.

\subsection{Reinforcement Learning for Imbalanced or Cost-Sensitive Detection}
Deep Q-networks~\cite{mnih2015human} overestimate action values, which Double DQN reduces by decoupling action selection and evaluation~\cite{vanhasselt2016double}; dueling heads separate state value from action advantage~\cite{wang2016dueling}. Lin et al. cast imbalanced classification as a sequential decision process solved by DQN with larger rewards for minority-class samples~\cite{lin2020imbalanced}, and Yang et al. combined dueling and double Q-learning with a custom reward for multi-class imbalanced clinical prediction~\cite{yang2024drlimbalanced}. E-DDQN applies the same ingredients to cybersickness. Safe-RL research shows that penalties shape how risk-averse a learned policy is~\cite{garcia2015safe}, but the relation between asymmetric rewards and classical cost-sensitive thresholding~\cite{elkan2001cost} is rarely made explicit, and RL detectors are seldom compared with cost-weighted supervised models under the same tuning and threshold rules.

\subsection{Gap Addressed}
Table~\ref{tab:gaps} summarizes the gaps this study addresses. The common thread is that a safety benefit can only be attributed to a detector design when the evaluation excludes leakage, the comparison is matched, and the decision rule is stated in terms of the costs it encodes.

\begin{table*}[t]
\centering
\caption{Gaps in the evaluation of asymmetric-reward cybersickness detectors and how this study addresses them.}
\label{tab:gaps}
\footnotesize
\begin{tabular}{p{0.21\textwidth}p{0.27\textwidth}p{0.44\textwidth}}
\toprule
Gap & Evidence & This study \\
\midrule
Sample-level splits of 1-Hz session data & Sample-level splits on the same gameplay data~\cite{kundu2022truvr}; leakage literature~\cite{kapoor2023leakage,kaufman2012leakage,roberts2017cv} & Quantified leakage; session-linkage groups, nested grouped evaluation, train-only preprocessing, protocol fixed before testing \\
Relation of asymmetric rewards to cost-sensitive learning left implicit & \cite{lin2020imbalanced,yang2024drlimbalanced} & Proof that the 5:1 reward equals the 5:1 cost-sensitive Bayes threshold (Eq.~\ref{eq:margin}); RL and cost-weighted models on one probability scale \\
RL detectors not compared with matched cost-sensitive or minimal baselines & -- & Seven models with equal budgets and rules; two exposure-time baselines; group-bootstrap inference with multiplicity adjustment \\
Robustness assessed with optimization-based input attacks only & \cite{kundu2025securing,ahmed2025adversarialvr} & Bounded, physically motivated telemetry perturbations with adaptive evasion and false-alarm worst cases \\
\bottomrule
\end{tabular}
\end{table*}
